Untuk setiap titik
\mathl{P \in M}{} kita pilih sebuah lingkungan peta terbuka
\mathl{P \in U_P}{} dengan sebuah peta
\maabbdisp {\alpha_P} {U_P } {V_P
} {}
dan
\mathl{V_P \subseteq \R^d}{.} Dengan beralih ke sebuah lingkungan bola dari
\mathl{\alpha_P(P) \in V_P}{} beserta prabayangnya, kita dapat mengandaikan bahwa $V_P$ adalah bola-bola terbuka dengan jari-jari paling besar
\mathl{{ \frac{   1 }{  3 } }}{}. Keluarga

\mathbed {U_P} {}
 {P\in M} {}
 {} {} {} {,}
menutupi manifold tersebut. Karena $M$ kompak, terdapat berhingga banyak titik
\mathl{P_1 , \ldots , P_n}{} sedemikian sehingga

\mathbed {U_i=U_{P_i}} {}
 {i=1 , \ldots , n} {}
 {} {} {} {,}
juga menutupi manifold tersebut. Kita menempatkan bola-bola terbuka
\mathl{V_i=V_{P_i}}{} di dalam
\zusatzklammer {\anfuehrung{$\R^d$ baru}{}} {} {}
dengan titik-titik pusat

\mathdisp {(1,0 , \ldots , 0), (2,0 , \ldots , 0) , \ldots , (n,0 , \ldots , 0)} { . }

Berdasarkan syarat jari-jari, bola-bola ini saling lepas. Kita meninjau himpunan terbuka terbatas
\mathl{U= \bigcup_{i=1}^n V_i}{.} Peta-peta $\alpha_i$ menghasilkan pemetaan-pemetaan kontinu

\mathdisp {\varphi_i :V_i \stackrel{ \left(  \alpha_i  \right)^{-1} }{\longrightarrow} U_i \longrightarrow M} { . }

Karena himpunan-himpunan tersebut saling lepas, dengan menetapkan
\mathl{\varphi(z)= \varphi_i(z)}{,} jika
\mathl{z \in V_i}{} berlaku, kita memperoleh sebuah pemetaan kontinu
\maabbdisp {\varphi} { U } {M
} {.}
Karena sifat tutupan tersebut, pemetaan ini surjektif.
